\chapter{Trader}\label{ch:in:trader}

On 4 May 2021 the largest US derivatives exchange group announced that it would not reopen the trading pits it had
closed during the pandemic, except one options pit; by 2025, 93\% of its volume was electronic. The trader's job had
moved from shouting to clicking, and from clicking to supervising algorithms that quote and trade faster than any person.
What remains a trader's is judgement: which risks to hold, how to set the machines, when to stop them, and how to answer
for them. This chapter opens Part~III, one chapter per role, and each follows the same plan: what the role is at each kind
of employer, how it has changed, what controls and rules bind it, a documented or composite day, the pay evidence of
chapter~14 for the role, and one analysis that shows what the role's work consists of. For the trader the analysis is a
question every desk faces: how many algorithms one person can watch.

\begin{center}\footnotesize
\begin{tabular}{@{}p{3.2cm}p{9.4cm}@{}}
\toprule
\multicolumn{2}{@{}l}{\textbf{Role card: trader}}\\
\midrule
employers & market makers, banks, hedge funds, platforms\\
horizon of decisions & seconds to days\\
P\&L & owns a book (\omterm{def:in:trader:kinds}{risk trader}) or none (\omterm{def:in:trader:kinds}{execution trader})\\
reports to & head of desk\\
occupation codes in filings & 13-2099.01, 13-2051, 41-3031; titles with ``trader'', ``trading analyst''\\
where the series teaches it & Book~1, chapter~2; Book~10, chapter~24; Book~11, chapter~27\\
\bottomrule
\end{tabular}
\end{center}

\section{Three traders: market maker, bank flow desk, hedge fund}

The word ``trader'' covers jobs that share little but a screen. Two definitions separate the main kinds.

\begin{definition}[Risk trader, execution trader]\label{def:in:trader:kinds}
A \emph{risk trader}\index{risk trader} decides which positions the firm holds, within limits, and owns the profit and
loss of a book. An \emph{execution trader}\index{execution trader} carries out orders decided by someone else, a portfolio
manager or a client, and is judged on the cost and quality of execution rather than on the position's profit.
\end{definition}

\begin{itemize}
\item \textbf{At a market maker} the trader is a \omterm{def:in:trader:kinds}{risk trader} whose book is the inventory the firm's quotes leave it with
(Book~11, chapters~1--3). One market maker describes its quantitative traders as ``in the thick of the action--closely
watching the markets, making quick decisions'' and adds that ``they build tools to automate trading decisions''; its
institutional traders ``interact with brokers and other trading teams''.
\item \textbf{At a bank's flow desk} the trader makes prices to the bank's clients through sales (Book~1, chapter~2) and
manages the resulting risk: a \omterm{def:in:trader:kinds}{risk trader} whose book is the franchise's flow (Book~9, chapter~24).
\item \textbf{At a hedge fund} the trader is usually an \omterm{def:in:trader:kinds}{execution trader} for the portfolio managers, working orders
through brokers and algorithms (Book~10, chapter~24); some funds also give traders risk of their own.
\end{itemize}
The three share the discipline of limits and the obligation to know, at every moment, what the firm holds; they differ in
whom they serve, what they own, and how their pay is judged (chapter~13).

\section{From click trading to supervising algorithms}

\begin{definition}[Click trader, algorithm supervisor]\label{def:in:trader:supervisor}
A \emph{click trader}\index{click trader} enters and manages orders by hand on an electronic screen. An
\emph{algorithm supervisor}\index{algorithm supervisor} is a trader in charge of one or more trading algorithms who sets
their parameters, monitors them in real time, handles their alerts and stops them when they misbehave.
\end{definition}

The pit trader of chapter~2 read the crowd; the \omterm{def:in:trader:supervisor}{click trader} read the screen; the \omterm{def:in:trader:supervisor}{algorithm supervisor} reads the
algorithms. The change was not only technical. EU rules for algorithmic trading make the supervisor a named duty: the
real-time monitoring of algorithmic trading ``shall be undertaken by the trader in charge of the trading algorithm or
algorithmic trading strategy, and by the risk management function'', and the staff in charge ``shall respond to
operational and regulatory issues in a timely manner''. A supervisor's day is therefore a queue: alerts arrive from the
algorithms, each needs attention, and an alert that waits is risk that runs unattended. How many algorithms one person
can supervise is a queueing question, and the chapter answers it in section~5.

\section{Controls and accountability: registrations, certification, kill switches}

\begin{definition}[Certification regime]\label{def:in:trader:certification}
A \emph{certification regime}\index{certification regime} is a rule that makes a firm assess, at hiring and at least
yearly, that people in named functions are fit and proper for them, and certify it; in the UK the functions include
proprietary traders who could cause significant harm and the people who approve or monitor trading algorithms.
\end{definition}

\begin{dated}{2026-07}{The rules a trader works under}\label{dat:in:trader:controls}
\textbf{US}: a securities trader at a broker-dealer who trades on its own account away from an exchange registers as a
Securities Trader Representative, after the Securities Industry Essentials and Series~57 examinations. \textbf{UK}: acting
as ``a proprietary trader whose activity involves, or might involve, a risk of significant harm to the firm or any of its
customers'' is a certification function, and so are approving the deployment of a trading algorithm and having
significant responsibility for monitoring whether an algorithm remains compliant (FCA Handbook, SYSC 27.8). \textbf{EU}: a
firm trading algorithmically ``shall be able to cancel immediately, as an emergency measure, any or all of its unexecuted
orders'' (the kill functionality), and must monitor its algorithms in real time through the trader in charge and its risk
function (Delegated Regulation 2017/589, Articles 12 and 16).
\end{dated}

The rules turn a trader's judgement into a documented responsibility. A kill switch (Book~11, chapter~27) exists at the
firm and at the venue; a loss limit stops a book when its losses reach a stated size; a pre-trade risk check rejects an
order that breaches a limit before it leaves the firm. The trader's part is to know which of them apply, to set the
limits the firm delegates, and to use the stop when it is needed rather than when it is too late. Chapter~24 describes the
risk and compliance functions on the other side of these controls.

\section{A documented day}

The rules and one firm's published description allow a composite day of an \omterm{def:in:trader:supervisor}{algorithm supervisor} at a market maker; it is
a composite, not any firm's schedule.
\begin{enumerate}
\item \textbf{Before the open}: check the algorithms' parameters, limits and kill paths; review overnight events and
positions; agree the day's plan with researchers and developers.
\item \textbf{The open}: the busiest minutes, when quotes widen and alerts come fastest; the supervisor watches fills,
positions and the algorithms' own health signals.
\item \textbf{The session}: adjust parameters as conditions change, answer alerts, investigate anything unexplained, and
escalate to risk and technology when an issue is beyond the desk.
\item \textbf{Incidents}: stop an algorithm, a product or the whole book when behaviour is not understood, and record why.
\item \textbf{After the close}: reconcile positions and profit and loss, review the day's alerts and fills, and feed what
was learned back into the models and the controls.
\end{enumerate}
The bank flow trader's day adds clients: prices requested through sales, trades to hedge, and a franchise to protect
(chapter~23). The hedge fund \omterm{def:in:trader:kinds}{execution trader}'s day follows the portfolio managers' orders and the brokers' algorithms.

\begin{dated}{2025-09}{What traders are paid: the public evidence}\label{dat:in:trader:pay}
US labour condition applications, fiscal 2025, role family ``trader'': median offered base \$195\,000 at market makers
(65 applications, 9 employers; interquartile range \$150\,000--250\,000) and \$235\,000 at banks (58 applications, 5
employers; \$152\,500--282\,500); cells for systematic funds and platforms are suppressed. The government's occupational
survey for securities, commodities and financial services sales agents in the securities industry, May 2025: median
\$103\,030, 10th to 90th percentile \$55\,130--309\,440. \omterm{def:in:how-pay-works:base}{Base salary} only in the first, most bonuses excluded from the
second (chapter~14).
\end{dated}

\begin{omfigure}
\begin{tikzpicture}
\begin{axis}[width=10.5cm,height=4.6cm,xmin=0,xmax=340,ymin=0.4,ymax=3.6,ytick=data,
  yticklabels from table={figdata/industry/16-trader/pay.csv}{source},yticklabel style={font=\scriptsize},
  xticklabel style={font=\footnotesize},xlabel={\small \$ thousand a year},grid=major,grid style={gray!20},
  table/col sep=comma]
\addplot[draw=none,mark=none,error bars/.cd,x dir=both,x explicit,error mark=none,error bar style={omDef,line width=1pt}]
  table[x=p50,y=pos,x error plus expr=\thisrow{p90}-\thisrow{p50},x error minus expr=\thisrow{p50}-\thisrow{p10}]
  {figdata/industry/16-trader/pay.csv};
\addplot[draw=none,mark=none,error bars/.cd,x dir=both,x explicit,error mark=none,error bar style={omDef!40,line width=7pt}]
  table[x=p50,y=pos,x error plus expr=\thisrow{p75}-\thisrow{p50},x error minus expr=\thisrow{p50}-\thisrow{p25}]
  {figdata/industry/16-trader/pay.csv};
\addplot[only marks,mark=|,mark size=5pt,line width=1.5pt,omThm] table[x=p50,y=pos]{figdata/industry/16-trader/pay.csv};
\end{axis}
\end{tikzpicture}

\omcaption{Traders' pay evidence: the 10th to 90th percentile (thin), the interquartile range (thick) and the median
(mark) of offered base in fiscal 2025 labour condition applications, beside the occupational survey's wages of securities
sales agents in the securities industry, May 2025. Data: \texttt{data/industry/lca\_ranges.csv},
\texttt{oews\_finance.csv}, through \texttt{in\_trader.pay}.}\label{fig:in:trader:pay}
\end{omfigure}

The filings show base salaries between \$150\,000 and \$250\,000 for most sponsored traders at both kinds of employer,
with the banks' median higher; the survey's occupation mixes traders with brokers and advisers, which is why its median
is half the filings'. Bonus, which chapter~14 found can multiply base several times at a market maker, is in neither.

\section{Tutorial: one supervisor, many algorithms}\label{tut:in:trader:tut}

\begin{tutorial}
\textbf{Goal.} Find how many algorithms one supervisor can watch while keeping alerts from waiting, under stated rates.
\textbf{End state:} \cref{fig:in:trader:pay,fig:in:trader:tail,fig:in:trader:capacity}.
\begin{tutsteps}
\item \textbf{Pay evidence.} \texttt{firm.roles.lca\_cells(rows, role)} with the role \texttt{trader} reads chapter~14's ranges for the role family,
with its suppressed cells; the occupational survey's row comes from \texttt{oews\_finance.csv}.
\item \textbf{The queue.} Each algorithm raises alerts at an illustrative 1 per hour, as a Poisson stream; handling takes
30 seconds on average; the standard is that at most 1\% of alerts wait more than one minute. With exponential handling
the waiting time exceeds $t$ with probability $\rho e^{-(\mu-\lambda)t}$ (Book~4, chapter~8).
\omcode{code/firm/roles/firm_roles.py}{79}{93}{The M/M/1 waiting-time tail and the largest number of algorithms one
supervisor can watch within the standard.}\label{lst:in:trader:mm1}
\item \textbf{A heavier tail.} \texttt{lindley\_tail} simulates the queue with lognormal handling times (mean 30 seconds,
standard deviation 45 seconds) by the Lindley recursion.
\end{tutsteps}
\end{tutorial}

With these rates the supervisor meets the standard for 7 algorithms: at 7 algorithms 0.89\% of alerts wait more than a
minute, and at 8, 1.03\%. The answer is sensitive to the alert rate: at half an alert an hour it is 15, at two alerts 3, at
four only 1 (\cref{fig:in:trader:capacity}). It is even more sensitive to the standard: allowing five minutes instead of
one raises the number to 71 at one alert an hour. And the tail of the handling time matters: with lognormal handling of
the same mean, 1.37\% of alerts wait more than a minute at 7 algorithms (\cref{fig:in:trader:tail}).

\begin{omfigure}
\begin{tikzpicture}
\begin{axis}[width=10.5cm,height=5.4cm,xmin=0.5,xmax=20.5,ymin=0,ymax=4.5,xlabel={\small algorithms supervised},
  ylabel={\small alerts waiting over a minute, \%},tick label style={font=\footnotesize},grid=major,grid style={gray!20},
  table/col sep=comma,legend style={legend cell align=left,font=\scriptsize,at={(0.03,0.97)},anchor=north west}]
\addplot[omDef,thick,mark=*,mark size=1.5pt] table[x=n,y=mm1]{figdata/industry/16-trader/tails.csv};
\addplot[omThm,thick,mark=square*,mark size=1.5pt] table[x=n,y=lognormal]{figdata/industry/16-trader/tails.csv};
\addplot[black!60,dashed] coordinates {(0.5,1) (20.5,1)};
\legend{exponential handling (closed form),lognormal handling (simulated),the 1\% standard}
\end{axis}
\end{tikzpicture}

\omcaption{Share of alerts that wait more than one minute against the number of algorithms, each raising one alert an
hour, with handling of mean 30 seconds. Rates are illustrative. Data: \texttt{in\_trader.tails} (400\,000 simulated alerts
per point).}\label{fig:in:trader:tail}
\end{omfigure}

\begin{omfigure}
\begin{tikzpicture}
\begin{axis}[width=10.5cm,height=5.2cm,ybar,area legend,bar width=12pt,ymode=log,log origin=infty,ymin=0.8,ymax=300,xmin=0.4,xmax=4.6,
  xtick={1,2,3,4},xticklabels={0.5,1,2,4},log ticks with fixed point,
  xlabel={\small alerts per algorithm per hour},ylabel={\small algorithms per supervisor},
  tick label style={font=\footnotesize},grid=major,grid style={gray!20},table/col sep=comma,
  legend style={legend cell align=left,font=\scriptsize,at={(0.5,-0.3)},anchor=north,/tikz/every even column/.append style={column sep=8pt}},
  legend columns=2,point meta=rawy,nodes near coords,nodes near coords style={font=\scriptsize}]
\addplot[fill=omDef!55,draw=omDef] table[x=pos,y=one_minute]{figdata/industry/16-trader/capacity.csv};
\addplot[fill=omThm!45,draw=omThm] table[x=pos,y=five_minutes]{figdata/industry/16-trader/capacity.csv};
\legend{1\% wait over one minute,1\% wait over five minutes}
\end{axis}
\end{tikzpicture}

\omcaption{How many algorithms one supervisor can watch, against the alert rate, for two response standards (exponential
handling of mean 30 seconds; log scale). Rates are illustrative. Data: \texttt{in\_trader.capacity}.}\label{fig:in:trader:capacity}
\end{omfigure}

\begin{example}[Pooling]\label{ex:in:trader:pool}
A desk runs 40 algorithms at one alert an hour each. Split into separate books of at most 7, it needs 6 supervisors to
meet the one-minute standard. With the same standard and one shared queue, 2 supervisors suffice: at 40 alerts an hour
only 0.17\% of alerts wait more than a minute, because an alert waits only when both are busy; two pooled supervisors
can take 76 algorithms and three 174. The price of pooling is that each supervisor must be able to handle any
algorithm's alert.
\end{example}

The lesson for a desk is not a number but a design: fewer, better alerts raise a supervisor's reach more than any other
change, which is why much of a supervising trader's work, and of the developers beside them (chapter~19), goes into
deciding what deserves an alert.

\textbf{What to change next.} Give alerts two priorities and serve the urgent ones first; let alert rates rise with market
volatility (chapter~12) and find the number of supervisors a desk needs on its busiest days.

\section{Build: the role registry and the supervision model}\label{bld:in:trader:roles}

\begin{build}
\bfield{Purpose} Record each role of Part~III as data, and hold one analysis per role; this first increment adds the
registry, the trader's card and the supervision model.
\bfield{Interface} \texttt{firm.roles}: \texttt{RoleCard(name, firm\_types, horizon, pnl, reports\_to, soc\_codes,
title\_patterns, books, chapter)}; \texttt{REGISTRY}, \texttt{register}, \texttt{card}; \texttt{lca\_cells(rows, role, fy,
level)}; \texttt{mm1\_wait\_tail(lam, mu, t)}; \texttt{max\_strategies(a, h, t, p)}; \texttt{lindley\_tail(lam, service, t,
n, rng)}. NumPy.
\bfield{Rules} A role is registered once; its pay cells come from chapter~14's table with suppression kept; the
supervision model's rates are the caller's, labelled illustrative.
\bfield{Acceptance tests} \texttt{code/firm/roles/tests/}: registering a role twice fails; the M/M/1 tail is 1 at or
above full load and matches the simulated queue within sampling error; the largest number of algorithms is monotone in
the alert rate and the standard.
\bfield{Stretch} Priorities and several supervisors (M/M/c); alerts that cluster in bursts; handling that depends on the
supervisor's load.
\end{build}

\begin{omsources}
\item CME Group, press release of 4 May 2021; Form 10-K for 2025.
\item FINRA, Series 57; FCA Handbook, SYSC 27.8; Commission Delegated Regulation (EU) 2017/589, Articles 12 and 16.
\item Optiver, ``Breaking down the trading industry''.
\item Chapter~14's tables from the Department of Labor's LCA files and the BLS occupational survey.
\end{omsources}

\section{Exercises}

\begin{exercise}[$\star$]\label{exo:in:trader:1}
Classify as \omterm{def:in:trader:kinds}{risk trader} or \omterm{def:in:trader:kinds}{execution trader}: a market maker's options trader; a hedge fund trader working a portfolio
manager's order; a bank's rates trader pricing a client's swap and hedging it.
\end{exercise}

\begin{exercise}[$\star$]\label{exo:in:trader:2}
In an M/M/1 queue with $\lambda=7$ alerts an hour and $\mu=120$, compute the probability that an alert waits more than one
minute.
\end{exercise}

\begin{exercise}[$\star$]\label{exo:in:trader:3}
Which three rules in the dated box would apply to an \omterm{def:in:trader:supervisor}{algorithm supervisor} at a UK firm trading on EU venues, and what does
each require of the person?
\end{exercise}

\begin{exercise}[$\star\star$]\label{exo:in:trader:4}
Show that at a fixed standard the number of algorithms scales roughly in inverse proportion to the alert rate when the
supervisor is lightly loaded.
\end{exercise}

\begin{exercise}[$\star\star$]\label{exo:in:trader:5}
Why is the occupational survey's median for securities sales agents about half the labour filings' median for traders?
\end{exercise}

\begin{exercise}[$\star\star$]\label{exo:in:trader:6}
With the same mean handling time, why does a lognormal handling time with a standard deviation of 45 seconds make waits
longer than an exponential one?
\end{exercise}

\begin{exercise}[$\star\star\star$]\label{exo:in:trader:7}
\emph{Coding.} With two supervisors sharing one queue (M/M/2), how many algorithms meet the one-minute standard at one
alert an hour? Simulate or use the Erlang C formula.
\end{exercise}

\begin{exercise}[$\star\star\star$]\label{exo:in:trader:8}
\emph{Find the flaw.} ``Our supervisors handle each alert in 30 seconds and there are 60 alerts an hour per person, so they
are only half busy and could watch twice as many algorithms.''
\end{exercise}

\section{Problem: One Trader, Many Algorithms}

\begin{problem}[{Weekend problem --- one trader, many algorithms}]\label{pb:in:trader:1}
A market maker is moving its options desk from click trading to supervised algorithms and asks how many algorithms each
trader should watch.

\medskip\noindent\textbf{Part I --- The role.}
\begin{enumerate}
\item Define \omterm{def:in:trader:kinds}{risk trader} and \omterm{def:in:trader:kinds}{execution trader}; give an example of each.
\item Describe the trader at a market maker, a bank flow desk and a hedge fund.
\item Define \omterm{def:in:trader:supervisor}{click trader} and \omterm{def:in:trader:supervisor}{algorithm supervisor}.
\item What did the largest US derivatives exchange group decide in 2021, and what share of its volume was electronic in
2025?
\item What do the EU's rules require of the trader in charge of an algorithm?
\end{enumerate}

\medskip\noindent\textbf{Part II --- Controls.}
\begin{enumerate}[resume]
\item Define a \omterm{def:in:trader:certification}{certification regime} and name the UK functions that concern traders.
\item What registration does a US securities trader at a broker-dealer need?
\item What is kill functionality, and who must have access to it?
\item Describe a composite day of an \omterm{def:in:trader:supervisor}{algorithm supervisor}.
\item What does the pay evidence say, and what does it leave out?
\end{enumerate}

\medskip\noindent\textbf{Part III --- The queue.}
\begin{enumerate}[resume]
\item State the model: arrivals, handling, standard.
\item Write the M/M/1 waiting-time tail and compute it at 7 and 8 algorithms.
\item How does the answer change with the alert rate?
\item How does it change with a five-minute standard?
\item How does it change with lognormal handling?
\end{enumerate}

\medskip\noindent\textbf{Part IV --- The verdict.}
\begin{enumerate}[resume]
\item State the \emph{named result}: the largest number of algorithms one supervisor can oversee with at most a 1\% chance
that an alert waits more than a minute, and its sensitivity to the alert rate.
\item Which change raises a supervisor's reach most?
\item Why is utilisation a poor guide to capacity here?
\item What should the desk measure before choosing the number?
\item In two sentences, answer the market maker.
\end{enumerate}
\end{problem}

\section{Interview questions}

\begin{interviewq}[$\star$ \iqroles{trader}]\label{iq:in:trader:1}
Your algorithm starts losing money in a way you do not understand. What do you do in the next sixty seconds?
\end{interviewq}

\begin{interviewq}[$\star$ \iqroles{trader, risk}]\label{iq:in:trader:2}
What is the difference between a kill switch and a loss limit?
\end{interviewq}

\begin{interviewq}[$\star\star$ \iqroles{trader, researcher}]\label{iq:in:trader:3}
Alerts arrive at 60 an hour and each takes 30 seconds. What fraction of the time is the supervisor busy, and what is the
mean wait in an M/M/1 queue?
\end{interviewq}

\begin{interviewq}[$\star\star$ \iqroles{developer, trader}]\label{iq:in:trader:4}
Design the alerts for an options market-making algorithm. Which would you page on, and which only log?
\end{interviewq}

\begin{interviewq}[$\star\star$ \iqroles{bank, trader}]\label{iq:in:trader:5}
A client asks for a price in size just before a data release. How do you decide the price and the hedge?
\end{interviewq}

\begin{interviewq}[$\star\star\star$ \iqroles{trader, researcher}]\label{iq:in:trader:6}
How would you decide whether a trader adds value to a set of algorithms, rather than only watching them?
\end{interviewq}
